% !PRINTED 635
\weberpage{635}
Since we are now in a position to determine the prime factors of $\psi_s(r)$ for every $s$, we proceed to decompose the number $(r,\eta)^m$ into its prime factors according to formula (11).

For this purpose we form, according to (15), the least positive residues $\varkappa$ of the numbers, $-n, -2n, \cdots -(m-2)n$ with respect to the modulus $m$, and find, apart from the order, the numbers
\[
\varkappa = 1, 2, \ldots, n-1, n+1, \ldots, m-1;
\]
the residues $0$ and $n$ do not occur among them, because $s$ receives no value for which $s$ or $s+1$ becomes divisible by $m$. Of these numbers, however, exactly $n-1$ are smaller than $n$, and so the factor $\mathfrak{p}_{n'}$ occurs that many times in the product $\psi_1\psi_2\ldots\psi_{m-2}$. Then $p$ contains this factor once more, and consequently, by (11), it occurs exactly $n$ times in $(r,\eta)^m$. Since no prime factors other than the $\mathfrak{p}_{n'}$ can divide $(r,\eta)^m$, the decomposition is hereby completely carried out:
\weq{16}{(r,\eta)^m = \prod^{n} \mathfrak{p}_{n'}^{\,n};}
here $n$ runs through the sequence of positive numbers that are smaller than $m$ and relatively prime to $m$, and from the congruence
\weq{17}{n n' \equiv 1\ (\mathrm{mod}\ m)}
$n'$ is in each case to be determined\footnote{This theorem is due to \gesperrt{Kummer}. Cf. Theorie der idealen Primfactoren der complexen Zahlen etc. Abhandlungen der Berliner Akademie, 1856.}.

Formula (16) holds for an even as well as for an odd $m$, and is still valid also in the case $m=2$, for which it gives a result already known, since then $(r,\eta) = (-1,\eta)$ agrees with a \gesperrt{Gauss} sum (Bd.\ I, \S.\,171).

But now, excluding the case $m=2$ from this point on, we shall give formula (16) a somewhat more general shape.

Let $m=q^{\varkappa}$ be an arbitrary prime power, only greater than $2$, and let $p$ be a prime different from $q$, with the property that $p-1$ is divisible by $q$ or, if $q=2$, by $4$.

Further, let $m_1$ be the greatest common divisor of $p-1$ and $m$, and
\[
m = m_1 m_2,
\]
%CONT>>
